\chapter{The Tile Renderer (Trainable Pathway)}\label{ch:tile-renderer}

\section{Overview}

\code{rmesh-trainable} implements the differentiable forward and backward
passes used by \code{rmesh-train} and the Python bindings. Both directions
are \emph{pure compute} --- no fixed-function rasterizer --- so the forward
intermediates (per-tet interval geometry, per-pixel transmittance state) are
exactly the quantities the backward kernel needs to replay. The screen is
divided into tiles (size 12 in training); a sorted list of
(tile, tet) pairs assigns each tile the tets that touch it in depth order
(built by \code{rmesh-tile} + the radix sort, see Chapter~\ref{ch:gpu-sorting}),
and one workgroup processes each tile.

Two forward implementations exist:
\begin{itemize}
\item \textbf{Interval tiled} (\code{interval_generate.wgsl} +
  \code{interval_tiled_rasterize.wgsl}) --- the current path, mirroring the
  interactive interval renderer's geometry.
\item \textbf{Software rasterizer baseline} (\code{rasterize_compute.wgsl})
  --- the pre-interval kernel that does ray--tet plane intersection
  in-shader; still used by \code{rmesh-train}'s loop and kept as a baseline.
\end{itemize}

\section{Forward stage 1 --- interval generation}

\code{interval_generate.wgsl} (\code{IntervalGeneratePipeline}) runs one
thread per visible tet, indexed through \code{compact_tet_ids}. It performs
the same silhouette decomposition as the interactive renderer
(Chapter~\ref{ch:interval-renderer}): Case 1 (interior vertex $\to$ 3-triangle
fan) or Case 2 (crossing edges $\to$ 4-triangle fan), writing 5 fixed vertex
slots per tet. Differences from the interactive variant:

\begin{itemize}
\item vertices are written in \emph{pixel} coordinates rather than NDC
  ($x_{\text{px}} = (x_{\text{ndc}}+1)\frac{W}{2}$,
  $y_{\text{px}} = (1-y_{\text{ndc}})\frac{H}{2}$), ready for tile scanline
  tests;
\item per-vertex records carry
  $(x, y, z^{\text{ndc}}_{\text{front}}, z^{\text{ndc}}_{\text{back}})$,
  the front/back color offsets, and front/back per-vertex normals;
\item \code{interval_meta[vis_idx]} packs the silhouette vertex count in the
  low 4 bits and the tet id above
  (\code{num_sil | (tet_id << 4)});
\item per-tet flat data \code{interval_tet_data[vis_idx]} $=
  (\text{base color}, \rho)$ with
  $\text{base color} = c + \nabla c \cdot (o_{\text{cam}} - v_0)$.
\end{itemize}

Behind-camera, out-of-frustum, and sub-pixel tets write degenerate slots and
\code{interval_meta = 0}.

\section{Forward stage 2 --- tiled rasterization}

\code{interval_tiled_rasterize.wgsl}
(\code{IntervalTiledRasterizePipeline}): one workgroup of 32 threads (one
warp) per tile, dispatched \code{(num_tiles, 1, 1)}. Workgroup shared memory
holds per-pixel state for up to $16\times16$ pixels: accumulated
premultiplied color with \emph{log-space transmittance} in \code{.w},
a normal+depth accumulator, and a 5-float Mip-NeRF-360 distortion-loss
state.

The tile's tets are traversed \textbf{front to back} by walking its
\code{tile_ranges} entry from the near end. Per tet:

\begin{enumerate}
\item Load the 5 silhouette vertices; lanes each own one pixel row and
  compute the scanline span $[x_l, x_r]$ of the silhouette polygon on that
  row, trimming pixels already saturated
  ($\log T < -13.8$, i.e.\ $T < 10^{-6}$); lane 0 builds a prefix sum of row
  widths.
\item All 32 lanes stride over the covered pixels. For each pixel: find the
  containing fan triangle, compute barycentrics, interpolate
  $z^{\text{ndc}}_{f}, z^{\text{ndc}}_{b}$ and the color offsets.
\item Linearize depths, apply the ray scale, and evaluate the shared volume
  integral: $\tau = \rho d$, $\varphi$, $w_0$, $w_1$,
  $c_{\text{premul}} = w_0 c_b + w_1 c_f$ (colors ReLU-clamped).
\item Composite front-to-back in log space:
  $T_j = e^{\log T}$; if $T_j \ge 10^{-6}$:
  $\text{color} \mathrel{+}= c_{\text{premul}} T_j$;
  $\log T \mathrel{-}= \tau$. Normal/depth and distortion accumulators are
  updated with the same $T_j$, and optional per-tet auxiliary channels
  accumulate $T_j\,\alpha\,\text{aux}[tet]$.
\end{enumerate}

Final per-pixel writes: \code{rendered_image} RGBA with
$\alpha = 1 - T_{\text{final}}$, a 4-channel \code{xyzd_image}
(normal + expected depth), the 5-channel distortion state, and aux channels.

The baseline \code{rasterize_compute.wgsl}
(\code{RasterizeComputePipeline}) has the same tile/warp/scanline structure
but intersects the pixel ray with the tet's four faces in-shader (slab
method over the shared \code{FACES} winding table) instead of reading
pre-decomposed intervals.

\section{Backward stage 1 --- per-tile reverse traversal}

\code{backward_interval_tiled.wgsl} (\code{BackwardIntervalTiledPipeline}),
one warp per tile, two bind groups (WebGPU 8-buffer limit): read-only inputs
(loss gradients \code{dl_d_image}, \code{dl_d_xyzd},
\code{dl_d_distortion}, \code{dl_d_aux}, the forward outputs, and the
interval buffers) and atomic accumulators
(\code{d_interval_verts}, \code{d_interval_tet_data}, \code{d_aux_data}).

The tile is traversed \textbf{back to front} --- the reverse of the forward
order --- and the forward per-pixel state is reconstructed by \emph{undo}
rather than storage: initialize from the final rendered pixel
($\log T = \log(1-\alpha)$), and per tet recompute the forward quantities,
then peel its contribution off:
\[
  \log T_{\text{prev}} = \min(\log T + \tau,\, 0), \qquad
  \text{color}_{\text{prev}} = \text{color} - c_{\text{premul}}\,
  T_{\text{prev}} .
\]
The chain rule then flows $dL/d(\text{image})$ into
$d\,c_{\text{premul}}$, $d\tau$ (via
$d w_0/d\tau = \varphi' + e^{-\tau}$, $d w_1/d\tau = -\varphi'$),
$d\rho$, $d d$, and $d z^{\text{ndc}}_{f/b}$ (through the depth
linearization Jacobian
$\partial z_{\text{view}} / \partial z_{\text{ndc}} =
n f (f-n) / (f - z_{\text{ndc}} (f-n))^2$), respecting every forward clamp
(optical-depth clamp, ReLU on colors, NDC range). Contributions from the
expected-depth and distortion outputs are handled the same way (distortion
undoes its moment accumulators). Pixel gradients are scattered to the three
fan-triangle vertices by barycentric weight.

\paragraph{Accumulation strategy.} Each lane keeps private accumulators;
a butterfly \code{subgroupShuffleXor} reduction sums the 32 lanes and
\emph{only lane 0} flushes with \code{atomicAdd} --- one atomic add per
channel per tet per tile instead of per pixel. Per-vertex gradient records
hold 10 floats
($dz_f, dz_b, d\text{off}_f, d\text{off}_b, d n_f, d n_b$); per-tet records
hold $(d\,\text{base color}, d\rho)$.

\section{Backward stage 2 --- chain back to parameters}

\code{interval_chain_back.wgsl} (\code{IntervalChainBackPipeline}), one
thread per visible tet, maps screen-space gradients back to mesh parameters.
It reconstructs the same Case-1/Case-2 topology and applies:

\begin{itemize}
\item \textbf{density}: pass-through into \code{d_densities[tet]};
\item \textbf{base color}: into \code{d_base_colors[tet]}; since
  $\text{base} = c + \nabla c \cdot (o_{\text{cam}} - v_0)$, also
  $d\nabla c \mathrel{+}= (o_{\text{cam}} - v_0)\, \Sigma d_{\text{bc}}$ and
  $d v_0 \mathrel{-}= \nabla c\, \Sigma d_{\text{bc}}$;
\item \textbf{NDC depth $\to$ vertices}:
  $\partial z_{\text{ndc}} / \partial v_{\text{world}} =
  (\text{VP}_{z\,\text{row}} - z_{\text{ndc}}\,\text{VP}_{w\,\text{row}}) /
  w_{\text{clip}}$, scattered by the barycentric / edge-mix weights of the
  case;
\item \textbf{color offsets $\to$ gradient and vertices}: from
  $\text{off} = \nabla c \cdot (v - o_{\text{cam}})$;
\item \textbf{normals}: silhouette and center vertex normal gradients
  scattered into \code{d_vertex_normals}.
\end{itemize}

All writes are \code{atomicAdd} --- per-vertex buffers are indexed by global
vertex id, so vertices shared between tets accumulate correctly.

\paragraph{Parameters receiving gradients:} vertex positions, per-tet
densities, per-tet base colors (the SH DC term --- the trainable path runs
with \code{sh_degree = 0}), per-tet color gradients, per-vertex normals, and
optional per-tet aux channels.

\section{Symbolic derivation and validation}

\code{scripts/derive_backward.py} derives every piece of this math in SymPy
and validates it with finite differences: the segment integral and its
derivatives, the log-space compositing undo, ray--plane intersection
Jacobians in compact cross-product form
($\partial n / \partial v_a = [v_b - v_c]_\times$, etc.), the color chain
with ReLU guards, and a full single-tet/single-pixel end-to-end chain.
GPU-side finite-difference gradient checks live in
\code{crates/rmesh-render/tests/gradient_test.rs} and
\code{crates/rmesh-trainable/tests/kernel_tests.rs}.

How this pathway is driven per training step --- tile-pair generation,
sorting, loss, and Adam --- is covered in Chapter~\ref{ch:training-and-python}.

